% ============================================================================
%  SPIE Methods, Secs. 2.3-2.8 -- narrative follows the order of discovery:
%
%    2.3  SAC: what we trained (equations, not prose)
%    2.4  Observed: recovery decays while success rises -> two modes
%    2.5  The modes are separated by mechanical state -- which the policy
%         cannot observe. This is why the standard objective cannot fix it.
%    2.6  Planned-path geometry corresponds to that mechanical state, giving
%         the policy a partial, deployable mode signal
%    2.7  Asymmetric SAC: critic values the modes correctly, policy acts on the
%         proxy -- the two halves of the fix
%    2.8  Evaluation protocol
%
%  Secs. 2.1-2.2 (anatomy generation; simulation, devices, observations,
%  reward) unchanged from spie_methods_asymmetric.tex and not repeated here.
%
%  Requires: \usepackage{amsmath, amssymb, booktabs}
%
%  ⚠ The correspondence probe at the end of Sec. 2.6 is NOT YET EXECUTED.
%    Eq. (6) is exact and verifiable from code; the CORRELATION with the
%    privileged fields is what remains to be measured.
%    See SPIE_METHODS.md inventory H1-H3.
% ============================================================================

\subsection{Reinforcement learning formulation}
\label{sec:rl}

Navigation is formulated as continuous control over the four device commands. The executed
command composes a learned residual with a scripted centerline follower,
%
\begin{equation}
a_t = \operatorname{clip}\!\big(a_{\mathrm{heur}}(s_t) + a_{\pi},\; a_{\min},\, a_{\max}\big),
\label{eq:residual}
\end{equation}
%
where $a_{\mathrm{heur}}$ is the scripted action and $a_{\pi}\sim\pi_{\phi}$ the learned residual,
so an untrained $\pi_{\phi}$ reproduces $a_{\mathrm{heur}}$ exactly while retaining authority to
override it anywhere within the actuation limits.

We train with soft actor-critic: two critics $Q_{\theta_1}, Q_{\theta_2}$ are fitted to the
bootstrapped return $y$, with the entropy term omitted from the backup, and the policy
$\pi_{\phi}$ maximizes the value regularized by an entropy temperature $\alpha$.
%
\begin{align}
\mathcal{L}_{Q}(\theta_i) &= \mathbb{E}_{(s,a,r,s')\sim\mathcal{D}}
   \Big[\big(Q_{\theta_i}(s,a) - y\big)^{2}\Big],
   \quad
   y = r + \gamma\,(1-d)\min_{j\in\{1,2\}} Q_{\bar{\theta}_j}(s',a'),
   \;\; a'\sim\pi_{\phi}(\cdot\mid s'),
   \label{eq:critic}\\[2pt]
\mathcal{L}_{\pi}(\phi) &= \mathbb{E}_{s\sim\mathcal{D},\, a\sim\pi_{\phi}}
   \Big[\alpha\log\pi_{\phi}(a\mid s) - \min_{j\in\{1,2\}} Q_{\theta_j}(s,a)\Big],
   \label{eq:actor}\\[2pt]
\mathcal{L}_{\alpha} &= \mathbb{E}\big[-\alpha\,\big(\log\pi_{\phi}(a\mid s) + \bar{\mathcal{H}}\big)\big].
   \label{eq:alpha}
\end{align}
%
Here $d$ is the terminal flag, $\gamma = 0.99$ the discount, and $\bar{\theta}_j$ the target
parameters, updated by Polyak averaging with $\tau = 0.005$. Both $\pi_{\phi}$ and $Q_{\theta_i}$
are two-layer perceptrons of width 256, with layer normalization in the critic bodies only. The
temperature $\alpha$ is auto-tuned by Eq.~\eqref{eq:alpha} against a target entropy
$\bar{\mathcal{H}} = +1.0$ rather than the conventional $-n_{\mathrm{actions}}$, with
$\log\alpha$ railed to $[-5,0]$, retaining exploration pressure late in training. Transitions are
drawn from the replay buffer $\mathcal{D}$ by prioritized sampling, with a fixed fraction of each
batch taken from episodes that reached the correct daughter branch. We call this
Configuration~A: $\pi_{\phi}$ and $Q_{\theta_i}$ receive the same deployable observation.

\subsection{Observed behavior: two navigation modes}
\label{sec:behavior}

Configuration~A improves steadily in episode success, but the trajectories show that this is not
the whole picture. Navigation resolves into two behaviors of very different frequency.
\emph{Centerline following} governs the large majority of steps and produces progress along the
route. \emph{Recovery}---retracting, reorienting and re-advancing after the device buckles or
enters a wrong branch---governs a small minority of steps, but they are the steps on which
episodes are lost.

We measure the second separately by detecting stalls from the step logs: a stall opens when the
guidewire is being advanced while progress along the route stays flat, and resolves when progress
passes the frontier at which it began. Stall resolution is a per-event rate and episode success a
per-episode rate, so the two are reported separately and never pooled.

Table~\ref{tab:recovery} shows the outcome. Across training, episode success rises while the
fraction of stalls resolved \emph{falls} by 24 percentage points; the two move in opposite
directions throughout. Configuration~A improves by entering stalls less often rather than by
escaping them better, and the aggregate metric hides this: monitoring episode success alone
records uninterrupted improvement while the ability to recover is being lost.

\subsection{The modes are separated by mechanical state, which the policy cannot observe}
\label{sec:why}

What distinguishes the two behaviors is not appearance but mechanics. In following, inserted
length converts into advancement and wall loading stays low and distributed. In a stall, insertion
continues while advancement does not: the applied length is absorbed into bending and contact, so
elastic and torsional energy accumulate and the load concentrates. The mode is therefore a
function of the mechanical state, and the simulator's contact forces, tensions and constraint
impulses separate the two cleanly.

The policy cannot see any of it. From a projection the device can be seen to have stopped
advancing---so the presence of a stall is partly observable---but not the state of that stall: how
much energy is stored, where the load is concentrated, or whether the contact is sticking or
slipping. Those are what determine \emph{which} action resolves it, and whether the correct
response is a millimetre of retraction, a large withdrawal, or rotation without retraction. Two
stalls that project identically can require opposite commands.

This is why the standard objective cannot resolve the competition, for two compounding reasons.
First, in the value estimate: $Q_{\theta_i}$ in Eq.~\eqref{eq:critic} regresses on the deployable
observation, so a recoverable configuration and a jammed one---indistinguishable in that
observation---are assigned a single averaged value. $Q_{\theta_i}$ cannot report that retraction is
worth taking in the first and wasted in the second, so the policy gradient through
Eq.~\eqref{eq:actor} receives nothing that separates them. Recovery is not merely under-rewarded;
it is unlearnable at the states where it matters.

Second, in the data: Eqs.~\eqref{eq:critic}--\eqref{eq:actor} are expectations over the state
distribution $d^{\pi_{\phi}}$ induced by the current policy. As following improves, stalled
configurations become rare under $d^{\pi_{\phi}}$ and their gradient contribution shrinks toward
zero.
Because episode success is dominated by the frequent behavior, an update that trades recovery for
following \emph{raises} the objective---the optimizer does not neglect the rare behavior, it
selects against it. Neither standard remedy applies: entropy regularization, Eq.~\eqref{eq:alpha},
diversifies actions at visited states but does not induce visitation of avoided ones, and replay
retains stalled transitions whose bootstrapped targets nonetheless drift under the increasingly
avoidant policy.

\subsection{Planned-path geometry corresponds to mechanical state}
\label{sec:correspondence}

% ⚠ PENDING -- the probe in the final paragraph is not yet implemented.
% Eq. (6) is an exact identity verifiable from code; that it CORRELATES with the
% privileged force fields is the claim still to be measured. Note also that any
% deployable feature numerically identical to a privileged dimension must be
% excluded from the probe, or it measures an identity rather than a correspondence.
% SPIE_METHODS.md inventory H1-H3.

Since the intervention is planned against a route, the device can be described in a reference
frame the operator already supplies---and in that frame part of the hidden mechanics becomes
visible. The relevant quantity is the discrepancy between inserted length and advancement,
%
\begin{equation}
\mathrm{slack}(t) \;=\; \ell_{\mathrm{ins}}(t) \;-\; s(t),
\label{eq:slack}
\end{equation}
%
where $\ell_{\mathrm{ins}}$ is inserted guidewire length and $s$ is arclength progress along the
planned route. Length inserted but not appearing as progress has been absorbed into bending, slack
loops and wall loading---exactly the accumulation that defines the stalled mode in
Sec.~\ref{sec:why}. Equation~\eqref{eq:slack} is a purely geometric scalar, computable from the
planned route and device tracking alone, that tracks the mechanical quantity the solver reports as
stored energy and contact loading. Related route-relative features---curvature ahead, cross-track
offset, heading error, distance to the next daughter branch---anticipate where loading and branch
selection will arise.

These features do not make the observation Markov. Frictional state is genuinely hysteretic, and no
instantaneous geometric quantity recovers whether a contact is sticking or slipping. What they
supply is a sufficient statistic for the component of the hidden state that is geometric in
origin---the stored-length component---which is the component that separates a recoverable
configuration from a jammed one. The planned route therefore converts part of an unobservable
mechanical problem into an observable geometric one, using information the clinical workflow
already provides.

We quantify how far this extends rather than assuming it. A probe predicts privileged stress
quantities---contact-impulse magnitude and distal nodal force---from route-relative features alone,
reported as area under the receiver operating characteristic curve for a binarized high-contact
state and as coefficient of determination for the continuous targets, each against a shuffled-label
control. Any deployable feature numerically identical to a privileged dimension is excluded, so the
probe measures correspondence rather than an identity.

\subsection{Asymmetric actor-critic}
\label{sec:asymmetric}

Section~\ref{sec:why} identified two failures, and they require different remedies. The policy's
blindness to mode is addressed by the route-relative features of
Sec.~\ref{sec:correspondence}, which enter its observation. The value estimate is addressed by
giving the critic the mechanical state directly:
%
\begin{equation}
Q_{\theta}\big(o^{\mathrm{priv}}, a\big),
\qquad
\pi_{\phi}\big(a \mid o^{\mathrm{dep}}\big),
\qquad
o^{\mathrm{dep}} \subset o^{\mathrm{priv}},
\label{eq:asym}
\end{equation}
%
with Eqs.~\eqref{eq:critic}--\eqref{eq:alpha} otherwise unchanged. The regression target of
$Q_{\theta_i}$ now depends on quantities that are Markov by construction---nodal forces and
velocities, contact impulse, accumulated twist---so a recoverable stall and a jammed one receive
\emph{different} values even when they project identically. The averaging that flattened the
recovery gradient in Sec.~\ref{sec:why} is removed at its source, and the policy gradient through
Eq.~\eqref{eq:actor} inherits a signal distinguishing a worthwhile retraction from a wasted one.

The two halves are complementary: the critic can now value mode-appropriate actions correctly, and
the policy has a deployable proxy by which to tell which mode it is in. Nothing privileged appears
at inference---the critics are discarded after training and only $\pi_{\phi}$ is executed---so the
policy is deployable by construction, with no gap between the trained and the deployed controller.
We call this Configuration~B; it differs from Configuration~A only in Eq.~\eqref{eq:asym}, which is
what licenses reading the difference between them as the effect of the asymmetry.

\subsection{Evaluation protocol}
\label{sec:eval}

Evaluation uses 98 fixed seeds across 50 procedurally generated anatomies, redrawn every two
episodes across 16 workers, with the clearance audit of Sec.~2.1 applied at generation time.
Actions are taken deterministically as the hyperbolic tangent of the policy mean, from the
deployable observation only. An episode succeeds if the guidewire tip comes within $5$\,mm
(three-dimensional Euclidean distance) of a target sampled at least $40$\,mm into the RCCA, within
600 control steps.

Anatomy identity is established by hashing branch-coordinate arrays, because the obvious
identifiers are unreliable: the planned path moves with the target, and the mesh fingerprint is
built from a seed reassigned on reset, so either can report a match between geometrically different
vessels.

Results are stratified by planned-path length into three bands, with cut points at $146$\,mm and
$210$\,mm, corresponding approximately to the common carotid artery, the mid-internal carotid
artery, and the cavernous segment. These are cut points on path length rather than anatomical
segmentation; the modelled vessel terminates at the internal carotid artery terminus, and the
middle cerebral artery is not represented.

Transfer is assessed on the patient-derived anatomy using its original segmented surface rather
than a reconstruction, with all other settings unchanged. Binomial (Wilson) confidence intervals
are reported; these understate the true uncertainty, since episodes are not fully independent and
the initial device twist is not seeded.
